% Part: second-order-logic
% Chapter: syntax-and-semantics
% Section: expressive-power

\documentclass[../../../include/open-logic-section]{subfiles}

\begin{document}

\olfileid{sol}{syn}{exp}
\olsection{Daya Ungkap}

\begin{explain}
Kuantifikasi tumrap variabel orde kapindho ndadekake daya ungkap
basa mundhak gedhe banget yen dibandhingake karo logika orde
kapisan. Kuantifikasi eksistensial orde kapindho ngidini kita
nyatakake anane fungsi utawa relasi kang nduweni sipat tartamtu.
Ing logika orde kapisan, siji-sijine cara kanggo iku yaiku
netepake simbol nonlogis (yaiku !!a{function} utawa
!!{predicate}) kanggo maksud kasebut. Kuantifikasi universal
orde kapindho ngidini kita nyatakake yen kabeh subhimpunan saka
domain, relasi ing domain, utawa fungsi saka !!{domain}
menyang !!{domain} nduweni sawijining sipat. Ing logika orde kapisan,
kita mung bisa nyatakake yen subhimpunan, relasi, utawa fungsi
kang diwenehake marang salah siji simbol nonlogis ing basa
nduweni sipat tartamtu. Nalika nyatakake anane subhimpunan,
relasi, utawa fungsi kang nduweni sawijining sipat, utawa yen
kabeh mau nduweni sipat iku, kita uga bisa nggunakake
kuantifikasi orde kapindho kanggo nemtokake sipat kasebut.
Iki ngidini kita netepake relasi kang ora bisa ditetepake ing
logika orde kapisan, lan nyatakake sipat domain kang ora bisa
dinyatakake ing logika orde kapisan.
\end{explain}

\begin{defn}
Yen $\Struct{M}$ iku !!a{structure} kanggo basa~$\Lang{L}$,
relasi~$R \subseteq \Domain{M}^2$ \emph{bisa didefinisikake}
ing~$\Lang{L}$ yen ana !!{formula}~$!A_R(x, y)$ kang mung
nduweni variabel bebas $x$ lan~$y$, saengga $R(a, b)$ bener
(yaiku, $\tuple{a,b} \in R$) yen lan mung yen
$\Sat{M}{!A_R(x, y)}[s]$ kanggo $s(x) = a$ lan $s(y) = b$.
\end{defn}

\begin{ex}
Ing logika orde kapisan kita bisa netepake relasi
identitas~$\Id{\Domain{M}}$ (yaiku
$\Setabs{\tuple{a,a}}{a \in \Domain{M}}$) nganggo formula
$\eq[x][y]$. Ing logika orde kapindho, relasi iki bisa
ditetepake \emph{tanpa~$\eq$}. Amarga yen $a$ lan $b$ iku
!!{element} kang padha saka~$\Domain{M}$, loro-lorone
kalebu !!{element}s saka subhimpunan-subhimpunan kang padha
saka~$\Domain{M}$ (amarga himpunan ditemtokake dening
!!{element}s kang ana ing njerone). Kosok baline, yen $a$ lan $b$ beda,
loro-lorone ora kalebu !!{element}s saka subhimpunan-subhimpunan
kang padha: upamane, $a \in \{a\}$ nanging $b \notin \{a\}$
yen $a \neq b$. Mula ``kalebu !!{element}s saka
subhimpunan-subhimpunan kang padha saka~$\Domain{M}$''
iku relasi kang bener kanggo $a$ lan $b$ yen lan mung yen
$a = b$. Relasi iki bisa dinyatakake ing logika orde kapindho,
amarga kita bisa nguantifikasi tumrap kabeh subhimpunan
saka~$\Domain{M}$. Mula !!{formula} ing ngisor iki
netepake $\Id{\Domain{M}}$:
\[
\lforall[X][(X(x) \liff X(y))]
\]
\end{ex}

\begin{prob}
Tuduhna yen $\lforall[X][(X(x) \lif X(y))]$ (elinga:
$\lif$, dudu~$\liff$!) netepake $\Id{\Domain{M}}$.
\end{prob}

\begin{ex}
Yen $R$ iku !!{predicate} rong-papan, $\Assign{R}{M}$ iku
relasi rong-papan ing~$\Domain{M}$. Sanajan rada mbingungake,
kita bakal nggunakake $R$ kanggo !!{predicate}~$R$ lan uga
kanggo relasi~$\Assign{R}{M}$ iku dhewe. \emph{Tutupan
transitif}~$R^*$ saka~$R$ iku relasi kang bener antarane
$a$ lan $b$ yen lan mung yen kanggo sawetara $c_1$, \dots,
$c_k$, pratelan $R(a,c_1)$, $R(c_1, c_2)$, \dots,
$R(c_k,b)$ bener. Iki kalebu kasus $k = 0$, yaiku yen
$R(a,b)$ bener, mula $R^*(a,b)$ uga bener. Mula
$R \subseteq R^*$. Sejatine, $R^*$ iku relasi paling cilik
kang ngemot~$R$ lan transitif. Ing logika orde kapindho kita
bisa nyatakake yen $X$ iku relasi transitif kang ngemot~$R$:
\begin{multline*}
  !B_R(X) \ident \lforall[x][\lforall[y][(R(x,y) \lif X(x, y))]] \land {}\\
\lforall[x][\lforall[y][\lforall[z][((X(x,y) \land X(y,z)) \lif X(x,
      z))]]].
\end{multline*}
Konjungsi kapisan nyatakake yen $R \subseteq X$, dene
konjungsi kapindho yen $X$ transitif.

Nyatakake yen $X$ iku relasi paling cilik kaya mangkono
tegese nyatakake yen relasi kasebut kalebu ing saben relasi kang ngemot
$R$ lan transitif. Mula kita bisa netepake tutupan transitif
saka~$R$ nganggo !!{formula}
\[
R^*(X) \ident !B_R(X) \land \lforall[Y][(!B_R(Y) \lif
  \lforall[x][\lforall[y][(X(x, y) \lif Y(x,y))]])].
\]
Kita nduweni $\Sat{M}{R^*(X)}[s]$ yen lan mung yen
$s(X) = R^*$. Tutupan transitif saka~$R$ ora bisa
dinyatakake ing logika orde kapisan.
\end{ex}

\end{document}
